\documentclass[12pt]{article}

% Layout
\usepackage[margin=1in]{geometry}
\usepackage{setspace}
\singlespacing

% Math
\usepackage{amsmath}
\usepackage{amssymb}

% Graphics and tables
\usepackage{graphicx}
\usepackage{booktabs}
\usepackage{threeparttable}
\usepackage{tabularx}
\usepackage{multirow}

% Citations (Chicago author-date)
\usepackage{chicago}

% Links
\usepackage[colorlinks=true,linkcolor=blue,citecolor=blue,urlcolor=blue]{hyperref}

% Title formatting
\usepackage{titling}
\pretitle{\begin{center}\large\bfseries}
\posttitle{\end{center}}
\preauthor{\begin{center}\normalsize}
\postauthor{\end{center}}
\predate{\begin{center}\normalsize}
\postdate{\end{center}}

% Tighter spacing for short version
\usepackage{enumitem}
\setlist{nosep}

% =====================================================================
% TITLE
% =====================================================================
\title{Information Permeability of Board Observer Networks\\A Research Proposal}

\author{
Harper Jung\thanks{Harvard University. Email: hjung@g.harvard.edu.}
}

\date{\today}

\begin{document}

\maketitle
\thispagestyle{empty}

\begin{abstract}
\noindent
Board observers (nonvoting attendees with full access to board meetings and materials but no fiduciary duty) are used by 82\% of venture capital entities, yet no empirical study has examined their role in corporate information flow. This proposal outlines a research design to test whether material events at private companies produce abnormal returns at connected public portfolio companies before public disclosure. I construct a person-level network linking 1,416 observers at private firms to their concurrent board directorships at public firms. Preliminary results suggest connected same-industry firms earn cumulative abnormal returns of 4.6 percentage points before M\&A announcements and 9.5 percentage points before bankruptcy filings, with no effect for routine events. Two regulatory shocks, the NVCA's 2020 removal of fiduciary language and the DOJ/FTC's January 2025 extension of Clayton Act Section~8 to observers, provide plausibly exogenous variation in opposite directions.

\medskip
\noindent \textbf{JEL Codes:} G30, G34, K22 \quad \textbf{Keywords:} board observers, information spillover, corporate governance, venture capital
\end{abstract}

\newpage
\setcounter{page}{1}

% =====================================================================
\section{Introduction}
\label{sec:intro}
% =====================================================================

Venture-backed governance operates through a three-tier architecture, namely directors with fiduciary duties and voting rights, board observers with full information access but no fiduciary duty or vote, and advisors with limited access. According to the NVCA's Q4 2023 survey, 82\% of VC entities appoint board observers, with universal adoption above \$500 million in AUM \cite{packin2025board}. Despite this prevalence, not a single empirical paper has studied observers. \citeN{carter2025} calls explicitly for such research, and \citeN{ewens2024board} acknowledge observers' existence in a footnote before excluding them.

This proposal asks whether the observer role creates information channels that would not exist under a traditional director-only board. Observers possess the same information as directors but face none of the legal constraints governing how directors handle that information, a disparity I term the \textit{accountability gap}. If this gap matters for information flow, the effects should be detectable in the stock prices of public companies connected to observers through concurrent board roles.

I construct a person-level network linking observers at private firms to their directorships at public firms. Using 57,000 material events at private companies from S\&P Capital IQ (2015--2025) and daily CRSP returns, I test whether connected public firms show abnormal returns before private firm events become public.

Three sets of preliminary findings motivate the proposed research. First, connected portfolio companies earn higher pre-event returns than non-connected stocks around the same events. Second, the effect is concentrated in material, confidential events, specifically M\&A announcements (+4.6 pp, $p=0.001$) and bankruptcy filings (+9.5 pp, $p=0.043$), and only when the connected firm operates in the same industry. Routine events produce no effect. Third, the channel responds to regulatory changes. The NVCA's 2020 removal of fiduciary language increased spillover ($p=0.001$), while the DOJ/FTC's January 2025 extension of Clayton Act Section~8 to observers decreased it ($p<0.001$).

The proposed contribution is threefold. First, I aim to document a private-to-public information channel operating through the Reg~FD exemption for private firms, distinct from the public-to-public spillovers studied in the interlocking directorates literature \cite{kang2008director,shi2025shared,fracassi2017corporate}. Second, I focus on pre-event leakage rather than post-event reputation spillovers. Third, I exploit two quasi-experimental regulatory shocks operating in opposite directions on the same mechanism. I frame the analysis as evidence of \textit{information permeability} rather than insider trading. The design can show that information reaches prices through the network, but does not identify the specific transmission mechanism.

% =====================================================================
\section{Institutional Background}
\label{sec:institutional}
% =====================================================================

\paragraph{Board observers in practice.} The NVCA model Investor Rights Agreement specifies that the company ``shall invite a representative of such Investor to attend all meetings of the Board of Directors in a nonvoting observer capacity and \ldots give such representative copies of all notices, minutes, consents, and other materials that it provides to its directors.'' In \textit{Obasi Investment Ltd.\ v.\ Tibet Pharmaceuticals} (3d Cir., 2019), the Third Circuit held that observers are not ``directors'' and do not owe fiduciary duties \cite{packin2025board}. Observer rights are contractual, arising from investor rights agreements rather than corporate charter.

\paragraph{The Reg~FD gap.} Regulation Fair Disclosure prohibits public companies from selectively disclosing material nonpublic information. Private companies are exempt. The private firms in my sample, where observers learn about pending acquisitions, bankruptcies, and strategic changes, are not subject to Reg~FD. The observer carries information from this unregulated environment to public companies where they serve as directors. The NVCA model IRA terminates observer rights ``immediately before the consummation of the IPO,'' implicitly acknowledging this incompatibility.

\paragraph{Two regulatory shocks.} In 2020, the NVCA removed the clause directing observers to ``act in a fiduciary manner with respect to all information so provided.'' After this change, the model template imposes no fiduciary obligation on observers, relying solely on confidentiality agreements. A textual analysis of 320 S-1 IRA exhibits shows that ``no fiduciary duty'' disclaimers increased from 46\% to 62\% after 2020 ($p=0.012$). In January 2025, the DOJ and FTC announced that board observers are subject to Clayton Act Section~8 interlocking directorate prohibitions, targeting precisely the same-industry connections I study. These two changes operate in opposite directions on the accountability gap (one loosening, one tightening), providing the proposed identification strategy.

% =====================================================================
\section{Data and Proposed Research Design}
\label{sec:data}
% =====================================================================

\subsection{Observer Network}

The network is constructed from three WRDS-linked sources. S\&P Capital IQ Professionals identifies 4,915 individuals with ``observer'' titles at 3,058 companies. For each, I retrieve concurrent public company positions from CIQ, supplemented with BoardEx (510,276 person matches via WRDS crosswalk) and Form~4 insider filings (790,551 matches). The resulting network contains 1,416 observers with 4,775 edges to public firms, 84\% representing director or chairman roles. Cross-referencing Form~4 filings reveals the CIQ-only network captures approximately half of all connections, motivating the supplementation. All results will be reported for both the original and supplemented networks.

\subsection{Events and Returns}

Events come from CIQ Key Developments (\texttt{wrds\_keydev}), totaling 400,886 events filtered to private companies, excluding CRSP-listed periods, earnings announcements, and conferences, yielding approximately 57,000 events. Key types include M\&A (Buyer/Target), bankruptcy, executive changes, and private placements. Daily returns from CRSP (2015--2025, 3.3 million observations) are used to compute market-adjusted CARs over windows from $[-30,-1]$ to $[0,+5]$, winsorized at the 1st/99th percentiles, excluding penny stocks.

\subsection{Empirical Strategy}

The baseline regression takes the following form.
\begin{equation}
\label{eq:baseline}
CAR_{i,e} = \beta_1 \cdot Connected_{i,e} + \beta_2 \cdot SameInd_{i,e} + \beta_3 \cdot Connected_{i,e} \times SameInd_{i,e} + \varepsilon_{i,e}
\end{equation}
where $i$ indexes portfolio stocks and $e$ indexes events. For each event, CARs are computed at all portfolio stocks, both connected and non-connected, so non-connected stocks serve as the control group. I report results across six clustering/FE specifications, with event-clustered standard errors as the most conservative.

For the regulatory shocks, I estimate on the connected-firm sample.
\begin{equation}
\label{eq:nvca}
CAR_{i,e} = \gamma_1 \cdot SameInd_{i,e} + \gamma_2 \cdot SameInd_{i,e} \times PostShock_e + \alpha_t + \varepsilon_{i,e}
\end{equation}
where $PostShock$ captures the NVCA 2020 or Clayton Act January 2025 break, with year FE and placebo tests at alternative dates.

\begin{table}[htbp]
\centering
\begin{threeparttable}
\caption{Summary Statistics}
\label{tab:summary}
\small
\begin{tabular}{lcc}
\toprule
 & Original Network & Supplemented Network \\
\midrule
\multicolumn{3}{l}{\textit{Panel A: Observer Network}} \\
\midrule
Unique observers & 1,140 & 1,416 \\
Observer-to-public-firm edges & 3,725 & 4,775 \\
Public portfolio companies & 1,694 & 2,744 \\
Private observed companies & 3,058 & 3,058 \\
Edges via director/chairman role & 83\% & 84\% \\
\midrule
\multicolumn{3}{l}{\textit{Panel B: Events (Filtered Sample)}} \\
\midrule
Total events (CIQ Key Dev) &  \multicolumn{2}{c}{$\approx$57,000} \\
\quad M\&A Buyer & 105 & 124 \\
\quad M\&A Target & 111 & 130 \\
\quad Bankruptcy & 128 & 146 \\
\quad Executive/Board Changes & 1,884 & 2,348 \\
\quad Private Placements & 4,344 & 5,102 \\
\quad Product/Client & 4,344 & 5,271 \\
Form D filings & \multicolumn{2}{c}{2,827} \\
\midrule
\multicolumn{3}{l}{\textit{Panel C: Returns}} \\
\midrule
CRSP daily observations & \multicolumn{2}{c}{3.3 million} \\
Portfolio stocks & 1,155 & 1,909 \\
Sample period & \multicolumn{2}{c}{Jan 2015 -- Dec 2025} \\
\bottomrule
\end{tabular}
\begin{tablenotes}
\footnotesize
\item \textit{Notes.} The original network is constructed from S\&P Capital IQ Professionals. The supplemented network adds BoardEx positions (via WRDS BoardEx--CIQ crosswalk) and Thomson Reuters Form~4 insider filings (via WRDS TR--CIQ crosswalk). Events are filtered to private companies, excluding CRSP-listed periods, earnings announcements, and conferences. CARs are market-adjusted, winsorized at the 1st/99th percentiles; stocks with average price below \$5 are excluded.
\end{tablenotes}
\end{threeparttable}
\end{table}

% =====================================================================
\section{Preliminary Results}
\label{sec:results}
% =====================================================================

\subsection{Baseline and Event-Type Heterogeneity}

\begin{table}[htbp]
\centering
\begin{threeparttable}
\caption{Baseline: Executive and Board Changes (CAR$[-5,-1]$)}
\label{tab:baseline}
\footnotesize
\begin{tabular}{lcccccc}
\toprule
 & (1) HC1 & (2) Event & (3) Stock & (4) Yr+Ev & (5) Yr+Stk & (6) Stk FE \\
\midrule
$Connected$ & 0.28** & 0.28** & 0.28** & 0.28** & 0.28** & 0.09 \\
 & (0.006) & (0.010) & (0.007) & (0.009) & (0.006) & (0.412) \\[0.2em]
$SameIndustry$ & $-$0.21 & $-$0.21 & $-$0.21 & $-$0.20 & $-$0.20 & $-$0.14 \\
 & (0.137) & (0.137) & (0.153) & (0.141) & (0.158) & (0.287) \\[0.2em]
$Conn \times SameInd$ & $-$0.05 & $-$0.05 & $-$0.05 & $-$0.06 & $-$0.06 & 0.11 \\
 & (0.796) & (0.796) & (0.815) & (0.773) & (0.793) & (0.604) \\
\midrule
$N$ & \multicolumn{6}{c}{253,782 (2,348 events)} \\
\bottomrule
\end{tabular}
\begin{tablenotes}
\scriptsize
\item \textit{Notes.} Market-adjusted CARs, supplemented network. Coefficients in pp. $p$-values in parentheses. $^{**}$ $p<0.05$.
\end{tablenotes}
\end{threeparttable}
\end{table}

For executive/board changes (2,348 events), connected firms earn 0.28 pp higher CAR$[-5,-1]$ than non-connected firms ($p=0.010$, event-clustered), surviving all specifications except stock FE (Table~\ref{tab:baseline}). Disaggregating by event type reveals the effect is concentrated in material events (Table~\ref{tab:eventtype}).

\begin{itemize}
\item \textbf{M\&A Buyer.} $Connected \times SameInd$ at CAR$[-30,-1]$ = +4.57 pp ($p=0.001$). Survives all non-stock-FE specifications at $p<0.01$. Neither $Connected$ alone ($-0.78$ pp, $p=0.318$) nor $SameInd$ alone ($-0.51$ pp, $p=0.708$) is significant. Only the interaction matters. The sample consists of 124 events and 28 connected same-industry observations.
\item \textbf{Bankruptcy.} $Connected \times SameInd$ at CAR$[-30,-1]$ = +9.45 pp ($p=0.043$). Consistent positive drift across all pre-event windows. Small sample (11 connected same-industry) makes estimates noisy.
\item \textbf{M\&A Target.} Effect appears only at CAR$[-1,0]$ = +1.08 pp ($p=0.027$). No pre-event drift, consistent with targets being the tightest board secrets.
\item \textbf{Routine events} (private placements, product announcements). No significant results.
\end{itemize}

\begin{table}[htbp]
\centering
\begin{threeparttable}
\caption{Event-Type Heterogeneity: $Connected \times SameIndustry$}
\label{tab:eventtype}
\small
\begin{tabular}{lccccc}
\toprule
 & CAR$[-30,-1]$ & CAR$[-10,-1]$ & CAR$[-5,-1]$ & CAR$[-1,0]$ & $N$ \\
\midrule
\multicolumn{6}{l}{\textit{Panel A: M\&A Buyer (observer's company acquiring)}} \\
\midrule
$Conn \times SameInd$ & 4.57*** & 2.43** & 1.31 & 0.42 & 14,226 \\
 & (0.001) & (0.039) & (0.185) & (0.327) & \\[0.2em]
$Connected$ & $-$0.78 & $-$0.36 & $-$0.15 & $-$0.08 & \\
 & (0.318) & (0.482) & (0.712) & (0.774) & \\[0.2em]
$SameIndustry$ & $-$0.51 & $-$0.22 & $-$0.04 & $-$0.11 & \\
 & (0.708) & (0.654) & (0.910) & (0.571) & \\[0.2em]
\quad Events & 124 &  &  &  & \\
\quad Connected same-ind & 28 &  &  &  & \\
\midrule
\multicolumn{6}{l}{\textit{Panel B: Bankruptcy}} \\
\midrule
$Conn \times SameInd$ & 9.45** & 6.28** & 4.05* & 1.22** & 16,412 \\
 & (0.043) & (0.047) & (0.074) & (0.033) & \\[0.2em]
$Connected$ & $-$0.31 & $-$0.18 & $-$0.22 & $-$0.09 & \\
 & (0.645) & (0.714) & (0.588) & (0.689) & \\[0.2em]
\quad Events & 146 &  &  &  & \\
\quad Connected same-ind & 11 &  &  &  & \\
\midrule
\multicolumn{6}{l}{\textit{Panel C: M\&A Target (observer's company being acquired)}} \\
\midrule
$Conn \times SameInd$ & 1.24 & 0.78 & 0.45 & 1.08** & 14,830 \\
 & (0.384) & (0.352) & (0.472) & (0.027) & \\[0.2em]
$Connected$ & $-$0.14 & 0.11 & 0.03 & $-$0.15 & \\
 & (0.802) & (0.774) & (0.921) & (0.512) & \\[0.2em]
\quad Events & 130 &  &  &  & \\
\quad Connected same-ind & 19 &  &  &  & \\
\midrule
\multicolumn{6}{l}{\textit{Panel D: Private Placements}} \\
\midrule
$Conn \times SameInd$ & 0.43 & 0.28 & 0.11 & 0.05 & 551,714 \\
 & (0.612) & (0.584) & (0.773) & (0.841) & \\[0.2em]
\quad Events & 5,102 &  &  &  & \\
\bottomrule
\end{tabular}
\begin{tablenotes}
\footnotesize
\item \textit{Notes.} This table reports the coefficient on $Connected \times SameIndustry$ from Equation~\ref{eq:baseline} estimated separately by event type and CAR window. All specifications use event-clustered standard errors and the supplemented network. Coefficients are in percentage points. $p$-values in parentheses. $^{*}$ $p<0.10$, $^{**}$ $p<0.05$, $^{***}$ $p<0.01$.
\end{tablenotes}
\end{threeparttable}
\end{table}

\subsection{Regulatory Shocks}

\begin{table}[htbp]
\centering
\begin{threeparttable}
\caption{Regulatory Shocks: NVCA 2020 and Clayton Act 2025}
\label{tab:shocks}
\small
\begin{tabular}{lcccc}
\toprule
 & CAR$[-30,-1]$ & CAR$[-10,-1]$ & CAR$[-5,-1]$ & CAR$[-1,0]$ \\
\midrule
\multicolumn{5}{l}{\textit{Panel A: NVCA 2020 --- Fiduciary Language Removal}} \\
\multicolumn{5}{l}{\textit{Sample: Connected firms, full period. Year FE, event-clustered SE.}} \\
\midrule
$SameInd \times Post2020$ & 2.86*** & 3.23*** & 1.67** & 0.54 \\
 & ($<$0.001) & (0.001) & (0.040) & (0.218) \\[0.2em]
$SameIndustry$ & $-$0.42 & $-$0.31 & $-$0.18 & $-$0.09 \\
 & (0.514) & (0.472) & (0.604) & (0.681) \\[0.3em]
Pre-2020 same-ind mean &  & 0.04\% &  & \\
 &  & (0.823) &  & \\
Post-2020 same-ind mean &  & 2.95\%*** &  & \\
 &  & ($<$0.001) &  & \\[0.3em]
Placebo: Jan 2024 break &  & 0.32 &  & \\
 &  & (0.510) &  & \\[0.2em]
Observations & \multicolumn{4}{c}{4,775} \\
\midrule
\multicolumn{5}{l}{\textit{Panel B: Clayton Act Jan 2025 --- Antitrust Extension to Observers}} \\
\multicolumn{5}{l}{\textit{Sample: Connected firms, post-2020. Year FE, event-clustered SE.}} \\
\midrule
$SameInd \times PostJan2025$ & $-$11.56*** & $-$10.19*** & $-$1.94* & $-$0.87 \\
 & (0.005) & ($<$0.001) & (0.088) & (0.214) \\[0.2em]
$SameIndustry$ & 1.83** & 2.41*** & 1.12** & 0.38 \\
 & (0.028) & ($<$0.001) & (0.041) & (0.305) \\[0.3em]
Placebo: Jan 2024 break &  & 1.61 &  & \\
 &  & (0.470) &  & \\[0.2em]
Observations & \multicolumn{4}{c}{3,218} \\
\midrule
\multicolumn{5}{l}{\textit{Panel C: Four-Period Trajectory (Same-Industry CAR$[-10,-1]$)}} \\
\midrule
Pre-2020 (baseline) & \multicolumn{4}{c}{0.04\% \quad ($p=0.823$)} \\
2020--Dec 2024 (NVCA loosened) & \multicolumn{4}{c}{2.95\%*** \quad ($p<0.001$)} \\
Jan--Sep 2025 (Clayton tightened) & \multicolumn{4}{c}{$-$2.56\%* \quad ($p=0.073$)} \\
Oct 2025+ (NVCA re-tightened) & \multicolumn{4}{c}{$-$2.04\% \quad ($p=0.134$)} \\
\bottomrule
\end{tabular}
\begin{tablenotes}
\footnotesize
\item \textit{Notes.} This table reports difference-in-differences estimates of regulatory shocks on same-industry information spillover. Panel~A estimates Equation~\ref{eq:nvca} on the connected-firm sample over the full period, with $Post2020=1$ for events in 2020 or later. Panel~B estimates the analogous specification on the post-2020 subsample, with $PostJan2025=1$ for events after January 2025. Panel~C reports mean same-industry CARs across four regulatory regimes. All specifications include year fixed effects. Coefficients are in percentage points. $p$-values in parentheses. $^{*}$ $p<0.10$, $^{**}$ $p<0.05$, $^{***}$ $p<0.01$.
\end{tablenotes}
\end{threeparttable}
\end{table}

\paragraph{NVCA 2020.} $SameInd \times Post2020$ at CAR$[-10,-1]$ = +3.23 pp ($p=0.001$). The pre-2020 same-industry CAR is 0.04 pp ($p=0.823$), while the post-2020 value is 2.95 pp ($p<0.001$). A placebo at January 2024 is null ($p=0.510$). Alternative break years confirm 2020 is the strongest.

\paragraph{Clayton Act 2025.} $SameInd \times PostJan2025$ at CAR$[-10,-1]$ = $-10.19$ pp ($p<0.001$). A placebo at January 2024 is null ($p=0.470$). The supplemented network yields $-4.88$ pp ($p=0.010$).

Two independent regulatory changes operating in opposite directions on the same mechanism, both with clean placebos. The post-January 2025 sample is still short. These estimates will become more precise as data accumulates.

\subsection{Preliminary Robustness}

The M\&A Buyer result survives SIC3 industry matching (4.94 pp, $p=0.004$), buy-and-hold abnormal returns (4.66 pp, $p=0.012$), and a continuous connection intensity measure (0.86 pp per link, $p=0.007$). Placebo tests (500 iterations) with random event dates reject the null at $p=0.048$ for M\&A Buyer and $p=0.022$ for Bankruptcy. Abnormal volume is null, and Form~4 insider trading patterns are suggestive but underpowered ($N=20$ pre-event same-industry trades).

% =====================================================================
\section{Expected Contributions and Next Steps}
\label{sec:conclusion}
% =====================================================================

If successful, this research would make three contributions. First, it would document a private-to-public information channel operating through the Reg~FD exemption, distinct from public-firm interlocking directorate spillovers. Second, it would identify pre-event information leakage rather than post-event spillovers. Third, it would provide quasi-experimental evidence from two opposite-direction regulatory shocks on the role of legal accountability in governing information flow.

\paragraph{Planned extensions.} (1)~Incorporating public observed company events to triple the sample. (2)~Calendar-time portfolio approach as alternative to event-time CARs. (3)~Additional post-shock data for Clayton Act and NVCA 2025 estimates. (4)~Network heterogeneity analysis by observer characteristics.

\paragraph{Channel identification.} A central extension is identifying the mechanism through which information reaches prices. I plan to investigate five candidate channels. (a)~\textit{Direct insider trading} by the observer, detectable through Form~4 trade timing (preliminary evidence is suggestive but underpowered). (b)~\textit{VC-level portfolio rebalancing}, where the observer's VC firm adjusts positions, testable via 13F holdings data. (c)~\textit{Strategic decision-making} at the connected public firm, where the observer-as-director influences corporate actions based on private board knowledge, testable through quarterly financials. (d)~\textit{Information diffusion} through the broader VC ecosystem to analysts, colleagues, and LPs, which predicts stronger effects for more central observers. (e)~\textit{Common investor rebalancing} by institutions holding stakes in both firms. The null result on abnormal volume tentatively argues against concentrated trading by a single party, but more granular tests are needed.

\paragraph{Limitations.} The network captures an estimated half of observer-to-public-firm connections. Same-industry samples for M\&A and bankruptcy are small (11--28 observations). The specific transmission mechanism is unidentified. The post-January 2025 sample is short.

% =====================================================================
% References
% =====================================================================
\newpage
\bibliographystyle{chicagoa}
\bibliography{../Bibliography_base}

\end{document}
